\section{The height carrier and the narrow crossover through zero}
\label{sec:crossover}

On a right-hand high-height chart use $L_z=\Log(z/(4\pi))$, continued
from the positive real axis, and define
\begin{align}
 A_o(t,z)&=\frac{\pi^{1/4}z^{7/4}}{16}
  \exp\left\{-\frac{\pi z}8+\frac t{16}(L_z^2-\pi^2/4)\right\},\\
 \Theta_o(t,z)&=\frac z4(L_z-1)+\frac{\pi tL_z}{16}+\frac{7\pi}8,
 \quad C_\pm=\tfrac12 A_o e^{\pm\ii\Theta_o},
 \quad \Eo=|C_+|+|C_-|.
 \label{eq:outer-arms}
\end{align}
The phase is unwrapped as written; each arm has half the cosine
amplitude. The envelope includes transverse growth. A real-axis
amplitude alone would not pay the buffered comparison.

Put $a(z)=\sqrt{z/(4\pi)}$, freeze $N=\lfloor a(X)\rfloor$, and set
\begin{align*}
 w_n(t,z)&=n^{-1/2}\exp\{t[(\log n)^2-L_z\log n]/4\},\\
 \varphi_n(t,z)&=\Theta_o(t,z)-(z/2+\pi t/8)\log n,\\
 \mathfrak C(v)&=\frac{\cos(2\pi(v^2-v)-\pi/8)}{\cos(2\pi v)},\\
 R_N(z)&=(-1)^N G(z)a(z)^{-1/2}\mathfrak C(a(z)-N),\\
 F_N(t,z)&=A_o(t,z)\sum_{n\le N}w_n(t,z)\cos\varphi_n(t,z)
                           +e^{-\pi^2t/64}R_N(z).
\end{align*}
Values at the apparent poles of $\mathfrak C$ are taken by removable
continuation. The identity
\begin{equation}
 \mathfrak C(v)+\mathfrak C(v-1)
        =2\cos(2\pi(v-1)^2-\pi/8)
 \label{eq:cutoff-identity}
\end{equation}
pays the adjacent cutoff change. This leading completion is simpler
than the high-order completion of Section~\ref{sec:negative}; it also
has a weaker error.

\begin{theorem}[Buffered signed crossover]
\label{thm:crossover}
On fixed buffered charts with $L_X=\log(X/(4\pi))$ and
$|t|L_X\le8$, let $t_-=\max(0,-t)$ and
$\delta=X^{-3/4}e^{t_-L_X^2/16}$. The completed canonical estimate is
\begin{align}
 |\partial_z^j(H-F_N)|&\le C_j\Eo\delta(1+L_X)^j,
                         &&0\le j\le3,\\
 |\partial_t(H-F_N)|&\le C\Eo\delta(1+L_X)^2,\\
 |(\partial_t+\partial_{zz})F_N|&\le C\Eo\delta(1+L_X).
 \label{eq:crossover-errors}
\end{align}
In the whole signed window $\delta=O(X^{-1/4})$; on its positive
half $\delta=X^{-3/4}$.
\end{theorem}
The proof retains both canonical arms, the complete arithmetic
population and cutoff payment on the larger strip. The time estimate
comes from the exact heat error identity and the actual residual; it
does not introduce a singular $1/t$ by differentiating an error
allowance. See \src{ansatz/notes/04_COMPLEX_STRIP_REMAINDER_AND_BUFFERED_CROSSOVER_20261025.md}{Ansatz Note 04}
and the supporting manuscript's
\src{ansatz/manuscript/main.tex}{buffered crossover estimates}
(\texttt{eq:buffered-crossover-jets}, \texttt{eq:buffered-crossover-time})
\cite{ansatz}.

This theorem explains a genuine but narrow two-sided passage through
zero. It does not extend the $e^{-L^2/4}$ recurrence estimate to zero.
An earlier paid chartwise negative-time join on $[-2,0]$ also uses
its own weaker representation and error account
(\src{ansatz/notes/06_NEGATIVE_BRIDGE_AND_TIME_JOIN_20261025.md}{Ansatz Note 06}).

\subsection{Retained arithmetic reaches the positive compact band}
The same frozen $F_N$, including its explicit completion, now has the
longer-band contract of Ansatz Note 26 \cite{retained}:
\begin{align}
 |\partial_z^j(H-F_N)|&\le C_j\Eo\Ls^j\epsilon_X,
       &&j\le3,\quad 4/L\le t\le1,\notag\\
 |(H-F_N)_t|&\le C_t\Eo\Ls^2\epsilon_X,
       &\epsilon_X&=\Ls^2X^{-3/4},\label{eq:retained-band}\\
 |\Res F_N|&\le C\Eo\epsilon_X,
 &|\partial_z\Res F_N|&\le C\Eo\Ls\epsilon_X.\notag
\end{align}
There are $2N=O(\sqrt X)$ arms. The proof supplies a growing-shift
initial-error lemma for the positive Gaussian's imaginary shifts
$O(tL)$. It uses a fixed inner strip and, outside it, the complete
ordinary Dirichlet tail, gamma ratios and floor-sensitive powers.
The full mode population costs $X^{1/4}$, while curvature/calibration
costs $\Ls^2/X$. Integrated Gaussian moments pay direct time before
using $1/t\le L/4$. These mechanisms explain the stated allowance;
a fixed-strip theorem alone would not cover the new shifts.

On $4\le tL\le8$ this is exactly the crossover formula, including
branches, normalization and completion. Its model difference is zero.
Consequently the same $F_N$ covers $0\le t\le1$ at the coarser
$\epsilon_X$ allowance, retaining the stronger inherited $X^{-3/4}$
allowance on $0\le tL\le8$. No time splice or derivative of $N$ is used.
The full proof and union are in the ansatz companion's
\src{ansatz/manuscript/08_positive_bridges.tex}{\texttt{thm:retained-positive-band}, \texttt{thm:positive-atlas}}.

Ansatz Note 24 separately proves $F_0=C_++C_-$ on $[1/2,1]$:
field/spatial $j\le3$ and direct-time allowances are respectively
$C_j\Eo\Ls^jX^{-t\log2/4}$ and
$C_t\Eo\Ls^2X^{-t\log2/4}$. The actual residual and its spatial
derivative are $C\Eo\Ls^2/X$ and $C\Eo\Ls^3/X$.
Its full-buffer calibration at $t=1/2$ is proved anew; a restricted-strip
reference theorem is not extended by citation \cite{small-positive}.
The ansatz \texttt{thm:small-positive-slab} contains the detailed proof.

Together with $F_1$ on $[3/4,33]$ and the compact saddle theorem on
$[32,T_*]$, these give a buffered high-height field atlas on $[0,T_*]$
for every fixed finite $T_*\ge33$, after the maximum of the actual
$T_*$-dependent onsets. Section~\ref{sec:positive} states the genuine
$F_1$ overlaps. This is an atlas of fields, with paid comparisons through
$H$, rather than one newly assembled smooth global field. Bounded-height
coverage, sharper physical consumers and useful joint cost are separate.
The stronger negative $e^{-L^2/4}$ precision still excludes zero; its
possible extension to $T\asymp L^{-2}$ is not supplied by this union.
